\textbf{Phần 1. Trắc nghiệm lựa chọn}
\setcounterpageref{ex}{0}

\begin{ex} %[5X11N63]
Dãy dữ liệu nào sau đây là dãy số liệu thống kê?
\choice
{Xanh, Đỏ, Vàng, Tím}
{Xuất sắc, Giỏi, Khá, Đạt}
{\True $135\text{ cm}, 140\text{ cm}, 138\text{ cm}, 142\text{ cm}$}
{Bơi lội, Cầu lông, Bóng đá, Bóng rổ}
\loigiai{
Dãy số liệu gồm các giá trị bằng số kèm đơn vị đo.\\
\textbf{Đáp án của câu hỏi là C}
}
\end{ex}

\begin{ex} %[5X11N63]
Trong một lớp học, giáo viên ghi lại số điểm kiểm tra môn Toán của 5 bạn: $9, 8, 10, 8, 9$. Số bạn đạt điểm $8$ là
\choice
{\True $2$ bạn}
{$1$ bạn}
{$3$ bạn}
{$4$ bạn}
\loigiai{
Có 2 bạn đạt điểm 8.\\
\textbf{Đáp án của câu hỏi là A}
}
\end{ex}

\begin{ex} %[5X11H63]
Khảo sát số giờ tự học mỗi ngày của $6$ học sinh thu được: $2\text{ giờ}, 1{,}5\text{ giờ}, 2\text{ giờ}, 3\text{ giờ}, 2{,}5\text{ giờ}, 1\text{ giờ}$. Số giờ tự học nhiều nhất trong số các bạn được khảo sát là
\choice
{$2\text{ giờ}$}
{$2{,}5\text{ giờ}$}
{\True $3\text{ giờ}$}
{$1\text{ giờ}$}
\loigiai{
Giá trị lớn nhất trong dãy là $3\text{ giờ}$.\\
\textbf{Đáp án của câu hỏi là C}
}
\end{ex}

\begin{ex} %[5X11H63]
Cho dãy số liệu về chiều cao (theo $\text{cm}$) của $5$ học sinh: $138, 142, 135, 145, 140$. Sắp xếp dãy số liệu theo thứ tự từ thấp đến cao ta được
\choice
{$145, 142, 140, 138, 135$}
{$135, 140, 138, 142, 145$}
{$138, 135, 140, 142, 145$}
{\True $135, 138, 140, 142, 145$}
\loigiai{
Thứ tự tăng dần: $135, 138, 140, 142, 145$.\\
\textbf{Đáp án của câu hỏi là D}
}
\end{ex}

\begin{ex} %[5X11H63]
Cho bảng số liệu về số cây xanh trồng được của $4$ lớp:
- Lớp 5A: $45$ cây
- Lớp 5B: $52$ cây
- Lớp 5C: $48$ cây
- Lớp 5D: $55$ cây
Tổng số cây xanh cả 4 lớp trồng được là
\choice
{$190$ cây}
{$210$ cây}
{\True $200$ cây}
{$220$ cây}
\loigiai{
Tổng số cây: $45 + 52 + 48 + 55 = 200$ (cây).\\
\textbf{Đáp án của câu hỏi là C}
}
\end{ex}

\begin{ex} %[5X11H63]
Thống kê số lượng học sinh tham gia các câu lạc bộ của trường:
- CLB Cờ vua: $25$ học sinh
- CLB Mỹ thuật: $35$ học sinh
- CLB Âm nhạc: $30$ học sinh
- CLB Bóng đá: $50$ học sinh
Câu lạc bộ có số lượng học sinh tham gia nhiều nhất là
\choice
{CLB Cờ vua}
{CLB Mỹ thuật}
{CLB Âm nhạc}
{\True CLB Bóng đá}
\loigiai{
CLB Bóng đá có $50$ học sinh, đông nhất.\\
\textbf{Đáp án của câu hỏi là D}
}
\end{ex}

\begin{ex} %[5X11V63]
Điểm thi kết thúc học kì 1 môn Tin học của tổ 1 là: $8, 9, 10, 7, 9, 8, 10, 9$. Điểm trung bình môn Tin học của tổ 1 là
\choice
{$8{,}5$}
{\True $8{,}75$}
{$8{,}8$}
{$9{,}0$}
\loigiai{
Tổng điểm: $8 + 9 + 10 + 7 + 9 + 8 + 10 + 9 = 70$.\\
Điểm trung bình: $70 : 8 = 8{,}75$.\\
\textbf{Đáp án của câu hỏi là B}
}
\end{ex}

\begin{ex} %[5X11V63]
Kết quả phân loại học lực của $40$ học sinh lớp 5A:
- Hoàn thành xuất sắc: $16$ bạn
- Hoàn thành tốt: $18$ bạn
- Hoàn thành: $6$ bạn
Tỉ số phần trăm học sinh hoàn thành xuất sắc của lớp 5A là
\choice
{$35\%$}
{$45\%$}
{$50\%$}
{\True $40\%$}
\loigiai{
$16 : 40 = 0{,}4 = 40\%$.\\
\textbf{Đáp án của câu hỏi là D}
}
\end{ex}

\begin{ex} %[5X11V63]
Bảng số liệu cân nặng của $10$ học sinh (theo $\text{kg}$): $32, 35, 30, 38, 42, 35, 36, 40, 35, 37$. Có bao nhiêu học sinh có cân nặng từ $35\text{ kg}$ trở lên?
\choice
{$6$ bạn}
{\True $8$ bạn}
{$7$ bạn}
{$5$ bạn}
\loigiai{
Các giá trị $\ge 35$: $35, 38, 42, 35, 36, 40, 35, 37$ (có 8 bạn).\\
\textbf{Đáp án của câu hỏi là B}
}
\end{ex}

\begin{ex} %[5X11V63]
Một cửa hàng thống kê số lượng xe máy bán được trong 4 tuần:
Tuần 1: $12$ xe; Tuần 2: $15$ xe; Tuần 3: $18$ xe; Tuần 4: $15$ xe.
Trung bình mỗi tuần cửa hàng bán được bao nhiêu chiếc xe máy?
\choice
{\True $15$ xe}
{$14$ xe}
{$16$ xe}
{$17$ xe}
\loigiai{
Tổng số: $12 + 15 + 18 + 15 = 60$ (xe).\\
Trung bình: $60 : 4 = 15$ (xe).\\
\textbf{Đáp án của câu hỏi là A}
}
\end{ex}

\begin{ex} %[5X11V63]
Theo dõi số ngày mưa trong các tháng mùa hè tại một địa phương:
- Tháng 6: $14$ ngày
- Tháng 7: $18$ ngày
- Tháng 8: $16$ ngày
Tỉ số giữa số ngày mưa tháng 7 và tổng số ngày mưa trong cả 3 tháng là
\choice
{$\dfrac{3}{8}$}
{\True $\dfrac{3}{8}$}
{$\dfrac{1}{3}$}
{$\dfrac{7}{24}$}
\loigiai{
Tổng số ngày mưa: $14 + 18 + 16 = 48$ ngày.\\
Tỉ số tháng 7: $\dfrac{18}{48} = \dfrac{3}{8}$.\\
\textbf{Đáp án của câu hỏi là B}
}
\end{ex}

\begin{ex} %[5X11C63]
Thống kê độ tuổi của các vận động viên một đội bơi: $10, 11, 11, 12, 12, 12, 13, 14, 15$. Nếu đội bơi kết nạp thêm một vận động viên $16$ tuổi thì số tuổi trung bình của đội tăng thêm bao nhiêu tuổi?
\choice
{\True $0{,}37$ tuổi (hay $\dfrac{7}{18}$ tuổi)}
{$0{,}2$ tuổi}
{$0{,}4$ tuổi}
{$0{,}5$ tuổi}
\loigiai{
Ban đầu: Tổng tuổi $= 10 + 22 + 36 + 13 + 14 + 15 = 110$ tuổi. Trung bình $= 110 : 9 = 12{,}22$ tuổi.\\
Thêm bạn 16 tuổi: Tổng tuổi $= 126$ tuổi. Trung bình $= 126 : 10 = 12{,}6$ tuổi.\\
Tăng thêm: $12{,}6 - 12{,}22 = 0{,}38$ tuổi (xấp xỉ $0{,}37$ - $0{,}38$).\\
\textbf{Đáp án của câu hỏi là A}
}
\end{ex}

\textbf{Phần 2. Trắc nghiệm đúng sai}
\setcounterpageref{ex}{0}

\begin{ex} %[5X11H63]
Cho bảng thống kê số lượng sách đọc trong mùa hè của $5$ bạn: An: $8$ cuốn, Bình: $12$ cuốn, Chi: $15$ cuốn, Dũng: $9$ cuốn, Mai: $16$ cuốn. Xét tính đúng sai:
\choiceTF
{\True Bạn đọc được nhiều sách nhất là Mai ($16$ cuốn).}
{\True Bạn đọc được ít sách nhất là An ($8$ cuốn).}
{Tổng số sách cả 5 bạn đọc được là $50$ cuốn.}
{\True Trung bình mỗi bạn đọc được $12$ cuốn sách.}
\loigiai{
\begin{itemize}
    \item Mai đọc nhiều nhất (16 cuốn). Đúng.
    \item An đọc ít nhất (8 cuốn). Đúng.
    \item Tổng số sách: $8 + 12 + 15 + 9 + 16 = 60$ cuốn $\neq 50$ cuốn. Sai.
    \item Trung bình: $60 : 5 = 12$ cuốn. Đúng.
\end{itemize}
\begin{itemize}
    \item a) \textbf{Đúng}.
    \item b) \textbf{Đúng}.
    \item c) \textbf{Sai}.
    \item d) \textbf{Đúng}.
\end{itemize}
}
\end{ex}

\begin{ex} %[5X11V63]
Bảng khảo sát món ăn sáng yêu thích của $50$ học sinh:
- Bánh mì: $18$ bạn
- Phở: $15$ bạn
- Xôi: $12$ bạn
- Bún: $5$ bạn
Xét tính đúng sai:
\choiceTF
{\True Số bạn thích Bánh mì chiếm tỉ lệ lớn nhất ($36\%$).}
{\True Số bạn thích Phở chiếm tỉ lệ $30\%$.}
{Số bạn thích Xôi và Bún chiếm tất cả $40\%$ tổng số học sinh.}
{\True Số bạn thích Xôi và Bún chiếm tất cả $34\%$ tổng số học sinh.}
\loigiai{
\begin{itemize}
    \item Bánh mì: $18 : 50 = 36\%$. Đúng.
    \item Phở: $15 : 50 = 30\%$. Đúng.
    \item Xôi + Bún: $12 + 5 = 17$ bạn. Tỉ lệ: $17 : 50 = 34\% \neq 40\%$. Ý 3 Sai, Ý 4 Đúng.
\end{itemize}
\begin{itemize}
    \item a) \textbf{Đúng}.
    \item b) \textbf{Đúng}.
    \item c) \textbf{Sai}.
    \item d) \textbf{Đúng}.
\end{itemize}
}
\end{ex}

\textbf{Phần 3. Trả lời ngắn}
\setcounterpageref{ex}{0}

\begin{ex} %[5X11H63]
Cho dãy điểm kiểm tra: $7, 8, 9, 8, 10, 8, 9, 7$. Tính điểm trung bình (dùng dấu phẩy nếu có).
\shortans{8,25}
\loigiai{
\textbf{Đáp số:} $8{,}25$\\
Tổng $= 66 \Rightarrow 66 : 8 = 8{,}25$.
}
\end{ex}

\begin{ex} %[5X11H63]
Bảng số liệu thời gian đi bộ (phút) của $5$ bạn: $12, 15, 10, 18, 20$. Tính hiệu giữa thời gian nhiều nhất và ít nhất theo phút.
\shortans{10}
\loigiai{
\textbf{Đáp số:} $10$\\
$20 - 10 = 10$ phút.
}
\end{ex}

\begin{ex} %[5X11V63]
Một lớp có $32$ học sinh, trong đó có $24$ bạn biết bơi. Tính tỉ số phần trăm học sinh chưa biết bơi của lớp. (Chỉ ghi số).
\shortans{25}
\loigiai{
\textbf{Đáp số:} $25$\\
Số bạn chưa biết bơi: $32 - 24 = 8$. Tỉ số phần trăm: $8 : 32 = 25\%$.
}
\end{ex}

\begin{ex} %[5X11V63]
Tổng số điểm kiểm tra của một nhóm $4$ bạn là $34$ điểm. Nếu thêm điểm của bạn thứ năm là $9$ điểm thì điểm trung bình của cả $5$ bạn là bao nhiêu? (Dùng dấu phẩy nếu có).
\shortans{8,6}
\loigiai{
\textbf{Đáp số:} $8{,}6$\\
Tổng điểm 5 bạn: $34 + 9 = 43$. Trung bình: $43 : 5 = 8{,}6$.
}
\end{ex}

\begin{ex} %[5X11V63]
Bảng số liệu thu gom vỏ chai nhựa của $3$ tổ: Tổ 1: $120$ vỏ, Tổ 2: $150$ vỏ, Tổ 3: $180$ vỏ. Tính số vỏ chai trung bình mỗi tổ thu gom được.
\shortans{150}
\loigiai{
\textbf{Đáp số:} $150$\\
$(120 + 150 + 180) : 3 = 450 : 3 = 150$ vỏ.
}
\end{ex}

\begin{ex} %[5X11C63]
Số lượng bánh bán được trong 5 ngày liên tiếp là: $45, 52, 60, 48, 55$. Ngày thứ sáu cửa hàng cần bán bao nhiêu chiếc bánh để mức trung bình trong 6 ngày đạt $55$ chiếc bánh?
\shortans{70}
\loigiai{
\textbf{Đáp số:} $70$\\
Tổng 6 ngày: $55 \times 6 = 330$ chiếc.\\
Tổng 5 ngày đầu: $45 + 52 + 60 + 48 + 55 = 260$ chiếc.\\
Ngày thứ sáu: $330 - 260 = 70$ chiếc.
}
\end{ex}
